\begin{lrachaptercard}{Axiomatic Equality --- Language, Structure, Theory}
\begin{multicols}{2}
\LraCardMathRow{Language}{\mathcal L_{=}:=\{=\}}
The pure language of equality has no constant symbols, no function symbols,
and no non-logical relation symbols. Its only distinguished symbol is equality,
read as identity.

\LraCardMathRow{Structure}{\mathcal M=(M,=^{\mathcal M})}
The domain \(M\) is nonempty. Equality is not freely interpreted:
\[
  =^{\mathcal M}=\Delta_M
  =\{(x,x):x\in M\}\subseteq M\times M.
\]

\columnbreak

\LraCardMathRow{Theory}{\mathrm{Eq}\subseteq\mathrm{Sent}(\mathcal L_=)}
The equality theory consists of the reflexivity axiom and Leibniz substitution
schema. Symmetry, transitivity, uniqueness, and well-definedness are derived
from these.

\LraCardMathRow{Satisfaction}{\mathcal M\models\mathrm{Eq}}
Every equality sentence holds when \(=\) is interpreted as actual identity on
the domain.
\end{multicols}
\end{lrachaptercard}

\begin{lratopiccard}{Axiomatic Equality --- Native Interface}
\begin{lrasignaturetable}
\LraSignatureRow{Carrier/context}{A nonempty domain \(M\); variables range over objects of \(M\).}
\LraSignatureRow{Native relation}{Equality \(x=y\), interpreted by the diagonal \(\Delta_M\subseteq M\times M\).}
\LraSignatureRow{Constants}{None in the pure equality language.}
\LraSignatureRow{Operations}{None in the pure equality language. Function symbols appear only after a richer first-order signature is supplied.}
\LraSignatureRow{Axioms}{Reflexivity and Leibniz substitution.}
\LraSignatureRow{Immediate theorems}{Symmetry, transitivity, and equality as an equivalence relation.}
\LraSignatureRow{Derived definitions}{Distinctness \(x\neq y\), uniqueness \(\exists!\), and finite-size constraints such as ``at least two'' and ``at most two.''}
\LraSignatureRow{Structural echo}{The diagonal \(\Delta_M:M\to M\times M\); equalizers detect when two maps have the same value.}
\end{lrasignaturetable}
\end{lratopiccard}

\begin{lraproofpatterncard}{Equality --- Proof Patterns}
\begin{lrasignaturetable}
\LraSignatureRow{Prove symmetry}{Assume \(x=y\). Use Leibniz with \(\phi(z)\equiv(z=x)\) and reflexivity \(x=x\).}
\LraSignatureRow{Prove transitivity}{Assume \(x=y\) and \(y=z\). Use Leibniz with \(\phi(w)\equiv(x=w)\).}
\LraSignatureRow{Use substitution}{Choose the property \(\phi\) whose truth should transfer across \(x=y\).}
\LraSignatureRow{Prove uniqueness}{Show existence of a witness, then show any two witnesses are equal.}
\LraSignatureRow{Prove well-definedness}{Show the proposed output does not change when equal representatives or equal inputs are substituted.}
\LraSignatureRow{Use congruence}{For functions, replace equal inputs inside \(f\); for relations, replace equal inputs inside \(R\).}
\end{lrasignaturetable}
\end{lraproofpatterncard}
